% Options for packages loaded elsewhere
\PassOptionsToPackage{unicode}{hyperref}
\PassOptionsToPackage{hyphens}{url}
\PassOptionsToPackage{dvipsnames,svgnames,x11names}{xcolor}
%
\documentclass[
  letterpaper,
  DIV=11,
  numbers=noendperiod]{scrartcl}

\usepackage{amsmath,amssymb}
\usepackage{lmodern}
\usepackage{iftex}
\ifPDFTeX
  \usepackage[T1]{fontenc}
  \usepackage[utf8]{inputenc}
  \usepackage{textcomp} % provide euro and other symbols
\else % if luatex or xetex
  \usepackage{unicode-math}
  \defaultfontfeatures{Scale=MatchLowercase}
  \defaultfontfeatures[\rmfamily]{Ligatures=TeX,Scale=1}
\fi
% Use upquote if available, for straight quotes in verbatim environments
\IfFileExists{upquote.sty}{\usepackage{upquote}}{}
\IfFileExists{microtype.sty}{% use microtype if available
  \usepackage[]{microtype}
  \UseMicrotypeSet[protrusion]{basicmath} % disable protrusion for tt fonts
}{}
\makeatletter
\@ifundefined{KOMAClassName}{% if non-KOMA class
  \IfFileExists{parskip.sty}{%
    \usepackage{parskip}
  }{% else
    \setlength{\parindent}{0pt}
    \setlength{\parskip}{6pt plus 2pt minus 1pt}}
}{% if KOMA class
  \KOMAoptions{parskip=half}}
\makeatother
\usepackage{xcolor}
\setlength{\emergencystretch}{3em} % prevent overfull lines
\setcounter{secnumdepth}{5}
% Make \paragraph and \subparagraph free-standing
\ifx\paragraph\undefined\else
  \let\oldparagraph\paragraph
  \renewcommand{\paragraph}[1]{\oldparagraph{#1}\mbox{}}
\fi
\ifx\subparagraph\undefined\else
  \let\oldsubparagraph\subparagraph
  \renewcommand{\subparagraph}[1]{\oldsubparagraph{#1}\mbox{}}
\fi


\providecommand{\tightlist}{%
  \setlength{\itemsep}{0pt}\setlength{\parskip}{0pt}}\usepackage{longtable,booktabs,array}
\usepackage{calc} % for calculating minipage widths
% Correct order of tables after \paragraph or \subparagraph
\usepackage{etoolbox}
\makeatletter
\patchcmd\longtable{\par}{\if@noskipsec\mbox{}\fi\par}{}{}
\makeatother
% Allow footnotes in longtable head/foot
\IfFileExists{footnotehyper.sty}{\usepackage{footnotehyper}}{\usepackage{footnote}}
\makesavenoteenv{longtable}
\usepackage{graphicx}
\makeatletter
\def\maxwidth{\ifdim\Gin@nat@width>\linewidth\linewidth\else\Gin@nat@width\fi}
\def\maxheight{\ifdim\Gin@nat@height>\textheight\textheight\else\Gin@nat@height\fi}
\makeatother
% Scale images if necessary, so that they will not overflow the page
% margins by default, and it is still possible to overwrite the defaults
% using explicit options in \includegraphics[width, height, ...]{}
\setkeys{Gin}{width=\maxwidth,height=\maxheight,keepaspectratio}
% Set default figure placement to htbp
\makeatletter
\def\fps@figure{htbp}
\makeatother
\newlength{\cslhangindent}
\setlength{\cslhangindent}{1.5em}
\newlength{\csllabelwidth}
\setlength{\csllabelwidth}{3em}
\newlength{\cslentryspacingunit} % times entry-spacing
\setlength{\cslentryspacingunit}{\parskip}
\newenvironment{CSLReferences}[2] % #1 hanging-ident, #2 entry spacing
 {% don't indent paragraphs
  \setlength{\parindent}{0pt}
  % turn on hanging indent if param 1 is 1
  \ifodd #1
  \let\oldpar\par
  \def\par{\hangindent=\cslhangindent\oldpar}
  \fi
  % set entry spacing
  \setlength{\parskip}{#2\cslentryspacingunit}
 }%
 {}
\usepackage{calc}
\newcommand{\CSLBlock}[1]{#1\hfill\break}
\newcommand{\CSLLeftMargin}[1]{\parbox[t]{\csllabelwidth}{#1}}
\newcommand{\CSLRightInline}[1]{\parbox[t]{\linewidth - \csllabelwidth}{#1}\break}
\newcommand{\CSLIndent}[1]{\hspace{\cslhangindent}#1}

\usepackage{booktabs}
\usepackage{longtable}
\usepackage{array}
\usepackage{multirow}
\usepackage{wrapfig}
\usepackage{float}
\usepackage{colortbl}
\usepackage{pdflscape}
\usepackage{tabu}
\usepackage{threeparttable}
\usepackage{threeparttablex}
\usepackage[normalem]{ulem}
\usepackage{makecell}
\usepackage{xcolor}
\KOMAoption{captions}{tableheading}
\makeatletter
\makeatother
\makeatletter
\makeatother
\makeatletter
\@ifpackageloaded{caption}{}{\usepackage{caption}}
\AtBeginDocument{%
\ifdefined\contentsname
  \renewcommand*\contentsname{Table of contents}
\else
  \newcommand\contentsname{Table of contents}
\fi
\ifdefined\listfigurename
  \renewcommand*\listfigurename{List of Figures}
\else
  \newcommand\listfigurename{List of Figures}
\fi
\ifdefined\listtablename
  \renewcommand*\listtablename{List of Tables}
\else
  \newcommand\listtablename{List of Tables}
\fi
\ifdefined\figurename
  \renewcommand*\figurename{Figure}
\else
  \newcommand\figurename{Figure}
\fi
\ifdefined\tablename
  \renewcommand*\tablename{Table}
\else
  \newcommand\tablename{Table}
\fi
}
\@ifpackageloaded{float}{}{\usepackage{float}}
\floatstyle{ruled}
\@ifundefined{c@chapter}{\newfloat{codelisting}{h}{lop}}{\newfloat{codelisting}{h}{lop}[chapter]}
\floatname{codelisting}{Listing}
\newcommand*\listoflistings{\listof{codelisting}{List of Listings}}
\makeatother
\makeatletter
\@ifpackageloaded{caption}{}{\usepackage{caption}}
\@ifpackageloaded{subcaption}{}{\usepackage{subcaption}}
\makeatother
\makeatletter
\@ifpackageloaded{tcolorbox}{}{\usepackage[many]{tcolorbox}}
\makeatother
\makeatletter
\@ifundefined{shadecolor}{\definecolor{shadecolor}{rgb}{.97, .97, .97}}
\makeatother
\makeatletter
\makeatother
\ifLuaTeX
  \usepackage{selnolig}  % disable illegal ligatures
\fi
\IfFileExists{bookmark.sty}{\usepackage{bookmark}}{\usepackage{hyperref}}
\IfFileExists{xurl.sty}{\usepackage{xurl}}{} % add URL line breaks if available
\urlstyle{same} % disable monospaced font for URLs
\hypersetup{
  pdftitle={Digitization of the Australian Parliamentary Debates (1901-2022)},
  pdfauthor={Lindsay Katz; Rohan Alexander},
  colorlinks=true,
  linkcolor={blue},
  filecolor={Maroon},
  citecolor={Blue},
  urlcolor={Blue},
  pdfcreator={LaTeX via pandoc}}

\title{Digitization of the Australian Parliamentary Debates
(1901-2022)\footnote{Code is available at:
  https://github.com/lindsaykatz/hansard-proj}}
\author{Lindsay Katz\footnote{University of Toronto,
  lindsay.katz@mail.utoronto.ca} \and Rohan Alexander\footnote{University
  of Toronto, rohan.alexander@utoronto.ca}}
\date{October 27, 2022}

\begin{document}
\maketitle
\begin{abstract}
Public knowledge of what is said in parliament is a key tenet of
democracy. Parliamentary dialogue is also critical for political science
research. In Australia, following the UK tradition, the written record
of what is said in parliament is known as Hansard. While the Australian
Hansard has always been publicly available, it has been difficult to use
it for the purpose of large-scale macro and micro-level text analysis
because it has not been available as a dataset of sufficient quality to
be credibly analysed with statistical models. Following the lead of the
LiPaD project which achieved this for Canada, our project aims to
provide a new, comprehensive, high-quality database that captures all
proceedings of the Australian parliamentary debates from 1901 to the
present day using Hansard. Our dataset will be publicly available and
linked to other datasets such as election results. The creation of this
dataset will enable the exploration of questions that are not currently
possible to explore, serving as a valuable resource for both researchers
and policymakers. This work will provide a thorough description of the
creation and computational underpinnings of this database, followed by a
discussion of some applications. To this point, we have completed this
for 2011 to 2022.
\end{abstract}
\ifdefined\Shaded\renewenvironment{Shaded}{\begin{tcolorbox}[borderline west={3pt}{0pt}{shadecolor}, interior hidden, sharp corners, breakable, frame hidden, enhanced, boxrule=0pt]}{\end{tcolorbox}}\fi

\hypertarget{sec-intro}{%
\section{Introduction}\label{sec-intro}}

The official written record of parliamentary debates, formally known as
Hansard, plays a fundamental role in capturing the history of political
proceedings and facilitating the exploration of valuable research
questions. Originating in the British parliament, the production of
Hansard became tradition in a number of Commonwealth countries such as
Canada, the United Kingdom, and Australia (Vice and Farrell 2017). Given
the content and magnitude of these records, they have significance,
particularly in the context of political science research. In the case
of Canada, a team of researchers at the University of Toronto have
digitized Hansard from 1901 to 2019, an endeavor called the Linked
Parliamentary Data (LiPaD) project (Beelen et al. 2017). Having a
digitized version of Hansard enables researchers to perform advanced
analyses on these records using text analysis tools and statistical
modelling. Following the lead of the LiPaD project, in this paper we
introduce a novel database for the Australian Hansard. This is composed
of individual datasets for each sitting day in the House of
Representatives, containing details on everything said in parliament in
a form which can be readily used by researchers. With the development of
tools for large-scale text analysis, this database will serve as a
valuable resource for efficiently studying and exploring political
behaviour in Australia over time.

The House of Representatives performs a number of crucial governmental
functions, such as creating new laws and overseeing government
expenditure (House of Representatives 2018, ch.~1). Sittings of the
House follow a predefined order of business, regulated by procedural
rules called standing orders (House of Representatives 2018, ch.~8). A
typical sitting day in the Chamber has a number of scheduled proceedings
including debates on government business, 90 second member statements,
and Question Time (House of Representatives 2018, ch.~8). The Federation
Chamber was created in 1994 as a subordinate debate venue of the
Chamber. This allows for better time management of House business as its
proceedings occur simultaneously with those of the Chamber (House of
Representatives 2018, ch.~21). Sittings in the Federation Chamber are
quite different to those of the Chamber in terms of their order of
business and scope of discussion. Business matters discussed in the
Federation Chamber are limited largely to intermediate stages of bill
development, and the business of private Members (House of
Representatives 2018, ch.~21). It is the recording and compilation of
all of these proceedings on which Hansard is based, and it is
essentially, but not entirely, verbatim.

A week or so after each sitting day, a transcript is available for
download from the official Parliament of Australia website in both
static PDF and extensible markup language (XML) form. The PDF is the
official release, which is converted to typed text using Optical
Character Recognition (OCR) technology (Sherratt 2016). This conversion
allows for the production of the XML formatted transcript, though these
are not always perfect conversions and are particularly flawed in the
early decades of Hansard (Sherratt 2016). The PDF imposes formatting
designed for humans to read with ease, whereas XML is designed for
consistency and machine legibility. The nature of XML enables us to more
easily use code to manipulate these records at scale, motivating our
choice to develop our database solely using the XML formatted files. In
cases where we were unsure on how to proceed with processing the XML, we
defer first to the PDF, and then to the video recording of the
proceeding if available.

The end goal for this work is to create a publicly available database of
parsed Hansard transcripts from 1901 to present in a form which can be
readily analyzed. More specifically, we aim to present each transcript
in an ordered form with details on each speaker such as their political
affiliation, the time stamp of their statement (where available), and
whether they interjected (i.e.~spoke out of turn). Thus far, this has
been done for Hansard from mid-2011 to present.

We begin this paper with a description of what datasets for Australian
Hansard already exist (Section~\ref{sec-existing}). Next will be a
detailed description of the contents of our database
(Section~\ref{sec-description}), followed by a thorough discussion of
how this database was created (Section~\ref{sec-create}).

\hypertarget{sec-existing}{%
\section{Existing Datasets}\label{sec-existing}}

At present, the Hansard format that is available on the Parliament of
Australia website is not accessible for large scale analysis. To this
point, various researchers have had to create their own databases in an
attempt to produce usable, complete Hansard data. For instance, Sherratt
(2016) created an online, easy to read
\href{http://historichansard.net}{database} of Hansard from 1901 to 1980
using the XML files. These data can be navigated by year, parliament,
people, and bills (Sherratt 2016). As well, Leslie (2021) has released
data on all Members of the Australian House of Representatives from 1945
to 2019, and is developing an \texttt{R} package for these data.
Further, Alexander and Hodgetts (2021) created the
\texttt{AustralianPoliticians} \texttt{R} package, which contains a
number of datasets related to the political and biographical information
of Australian federal politicians who were active between 1901 and 2021.

A number of papers exist which utilize components of Australian Hansard
to explore various research topics. For example, Salisbury (2011) used
online Hansard records to explore the occurrences of unparliamentary
comments by members of the House, where The Speaker tells that member to
withdraw their remark. Rasiah (2010) worked with Question Time data from
Hansard transcripts during February and March of 2003, to investigate
resistance of politicians in answering questions about Iraq. In their
work, Fraussen, Graham, and Halpin (2018) use Hansard to quantify
political prominence by investigating strategic mentions of interest
groups by elected officials. Finally, the work of Alexander and
Alexander (2021) offers a dataset of the Australian Hansard, along with
an analysis of the effect of elections and changes in Prime Ministers
upon topics mentioned in parliament. Alexander and Alexander (2021)
created this database with the static PDF versions Hansard, using OCR to
digitize these files into text which is suitable for analysis. This
means there are considerable digitization errors especially in the first
half of the dataset.

While there is evidently a growing body of literature on this topic,
there is still no comprehensive database for Australian Hansard based on
XML that spans from 1901 to the present day. Our work serves to bridge
this gap.

\hypertarget{sec-description}{%
\section{Description of the Database}\label{sec-description}}

Our current database contains 670 comma-separated values (CSV) files,
one for each sitting day of the House of Representatives from 10 May
2011 to 08 September 2022. We used two scripts to produce these files,
which only have one subtle difference. Since the Federation Chamber was
renamed in mid-2012 from the original title of Main Committee, the
element under which all the Federation Chamber contents are nested in
the XML was renamed as well (House of Representatives 2018, ch.~21).
This difference in element name means that our original code for parsing
Federation Chamber content would not work, as the specified node path
would be incorrect. We accounted for this by creating a separate script
with this name change. Our first script parses Hansard from 14 August
2012 to 08 September 2022, that being everything after the Federation
Chamber name change, and the second script parses Hansard from 10 May
2011 to 28 June 2012. All output files in our database are formatted and
processed in the exact same way. Figure~\ref{fig-main15rows} provides a
snapshot of one of these files.

\begin{figure}

{\centering \includegraphics{main_ex15.png}

}

\caption{\label{fig-main15rows}First 15 rows of CSV file for Hansard on
25 February 2020}

\end{figure}

Each row of the CSV file contains an individual statement, with the
unique name ID, electorate, and party information of the speaker. Of
course, for general statements such as ``Honourable members'', these
variables cannot be specified. Each proceeding in the Chamber begins
with the business start, where the time that the Speaker took the chair
is documented, as well as the acknowledgment of the country and reading
of the prayers as governed by the standing orders of the House
(Department of the House of Representatives 2019, ch.~6). The Federation
Chamber also begins with a business start, where it is stated at what
time 5he Deputy Speaker takes the chair. For further clarification on
each variable in the resulting CSV, Table~\ref{tbl-vars} provides an
overview of each variable in the database.

\hypertarget{tbl-vars}{}
\begin{table}[H]
\caption{\label{tbl-vars}Summary and description of variables in our database }\tabularnewline

\centering
\begin{tabular}{ll}
\toprule
Variable & Description\\
\midrule
name & Name of speaker\\
order & Row number\\
speech\_no & Speech number\\
page.no & Page number statement can be found on in official Hansard\\
time.stamp & Time of statement\\
\addlinespace
name.id & Unique member identification code\\
electorate & Speaking member's electorate\\
party & Speaking member's party\\
body & Statement text\\
fedchamb\_flag & Flag for Federation Chamber (1 if Federation Chamber, 0 if Chamber)\\
\addlinespace
sub1\_flag & Flag for sub-debate 1 contents (1 if sub-debate 1, 0 otherwise)\\
sub2\_flag & Flag for sub-debate 2 contents (1 if sub-debate 2, 0 otherwise)\\
question & Flag for question (1 if question, 0 otherwise)\\
answer & Flag for answer (1 if answer, 0 otherwise)\\
q\_in\_writing & Flag for question in writing (1 if question in writing, 0 otherwise)\\
\addlinespace
interject & Flag for interjection (1 if statement is an interjection, 0 otherwise)\\
\bottomrule
\end{tabular}
\end{table}

The \texttt{name}, \texttt{page.no}, \texttt{time.stamp},
\texttt{name.id}, \texttt{electorate}, \texttt{party}, and \texttt{body}
variables all came directly from the XML contents. In addition to these
variables, we added a number of flags to enable easy filtering of
statements. For example, adding the \texttt{fedchamb\_flag} provides a
clear distinction between the proceedings of the Chamber with those of
the Federation Chamber. As well, the \texttt{sub1\_flag} and
\texttt{sub2\_flag} variables allow us to keep track of where various
statements are being parsed from in the XML document. Further details on
this XML structure will follow in Section~\ref{sec-create}. The
\texttt{question}, \texttt{answer}, and \texttt{q\_in\_writing} flags
were added to identify statements belonging to Question Time, and the
nature of these statements.\footnote{More on this in
  Section~\ref{sec-qa}} We also added a flag for interjections, which
are defined as statements made during one members speech by another
member who is not the Speaker or Deputy Speaker. Further, the
\texttt{speech\_no} variable allows us to keep track of the speech
number that each statement and interjection belongs to. For example, in
Figure~\ref{fig-main15rows}, statements belonging to speech number 3
spans five rows, including two interjections. Having the speech number
variable offers an easy way to group statements by speech or isolate
specific speeches of interest. Finally, the \texttt{order} variable was
added to keep track of the order in which everything was said, and was
created after all statements and interjections were separated onto their
own rows.

\hypertarget{sec-stats}{%
\subsection{Descriptive Statistics}\label{sec-stats}}

A visual exploration of our database offers interesting insight into
general characteristics of Hansard. First, we looked at the total number
of sitting days in the House of Representatives for each year from 2000
to 2021. Figure~\ref{fig-plot1} illustrates that the number of House
sitting days tends to fluctuate over time, where within this time span
the years 2013 and 2019 had the fewest sitting days (45), and 2014 had
the most sitting days (76).

\begin{figure}

{\centering \includegraphics[width=0.8\textwidth,height=\textheight]{paper_files/figure-pdf/fig-plot1-1.pdf}

}

\caption{\label{fig-plot1}Total number of sitting days per year
(2000-2021)}

\end{figure}

In Figure~\ref{fig-plot2}, we consider the average number of rows per
sitting day per year, disaggregated by Chamber and Federation Chamber
proceedings. We can see that 2014 has the highest average number of rows
per day, and 2021 has the lowest average number of rows per day. There
appears to be less fluctuation in the average number of rows per day in
the Federation Chamber than there is in that of the Chamber.

\begin{figure}

{\centering \includegraphics[width=0.8\textwidth,height=\textheight]{paper_files/figure-pdf/fig-plot2-1.pdf}

}

\caption{\label{fig-plot2}Average number of rows in Hansard per year, by
Chamber and Federation Chamber}

\end{figure}

Next we were interested in the percentage of rows in each Hansard which
are flagged as interjections, using our definition of interjections
mentioned above. We found that on average, approximately 18\% of rows
are characterized as interjections. Further, of all sitting days in our
database, 25 June 2014 had the greatest percentage of rows being
interjections, which was approximately 33\%. The day with the smallest
percentage of interjections was 23 March 2015, at approximately 3\%.
Based on Figure~\ref{fig-plot3}, there do not appear to be any
significant fluctuations or patterns in the percentage of rows that are
interjections across our time span.

\begin{figure}

{\centering \includegraphics[width=0.8\textwidth,height=\textheight]{paper_files/figure-pdf/fig-plot3-1.pdf}

}

\caption{\label{fig-plot3}Percentage of rows in Hansard that are
interjections, 10 May 2011 - 08 September 2022}

\end{figure}

Finally, Figure~\ref{fig-unique_plots} summarizes the number of unique
speakers, electorates, and parties found in Hansard over time. In all
three panels, the data exhibit the same general pattern over time for
both the House of Representatives and the Federation Chamber, where the
magnitude for the House data is consistently greater than that of the
Federation Chamber. Panels A and B in Figure~\ref{fig-unique_plots}
illustrate that the number of unique names and electorates follow a very
consistent trend across the entire time frame. Panel C shows that the
number of unique parties has increased slightly since 2018.

\begin{figure}

{\centering \includegraphics{unique_plots.png}

}

\caption{\label{fig-unique_plots}Number of unique names (A), electorates
(B), and parties (C) found on each sitting day.}

\end{figure}

\begin{table}

\end{table}

\hypertarget{sec-create}{%
\section{Creation of the Database}\label{sec-create}}

The approach to parsing contents of an XML document is heavily dependent
on its tree structure. As such, to create this database, we started by
looking at a single Hansard XML transcript. Doing so enabled us to
identify the various components of interest in the document, and how
each one can be parsed according to its corresponding structural form.
For example, the beginning of a speech may be followed by various
interjections and continuations, all of which are structured as unique
sub-elements containing further nested sub-elements within them. Parsing
was performed in \texttt{R} using the \texttt{XML} package. Focusing on
one transcript also allowed us to ensure that all key components of the
transcript were being parsed and captured with as much detail as
possible. Further, since we are looking at such a wide time span of
documents, there are expectedly a number of changes in the way they are
formatted. Some changes are as subtle as a differently named child-node,
while others are as extensive as a completely different nesting
structure. Smaller changes were accounted for as we became aware of
them, and embedded into the final scripts in a way that would not cause
issues for parsing more current Hansards with slightly different
formatting. As mentioned in Section~\ref{sec-intro}, the significant
changes in XML structure of Hansard preceding 2011 necessitates much
further development of our current scripts, which we are actively
working to complete.

Figure~\ref{fig-xml1} provides an example of the beginning of an XML
file for Hansard. The XML structure begins with a root or parent element
\texttt{"hansard"} (highlighted in blue), followed by a child element
\texttt{"session.header"} (highlighted in yellow) with sub-child
elements with information such as the date, and parliament number, which
are all highlighted in pink. Next there is a child element containing
everything that takes place in the Chamber, \texttt{"chamber.xscript"},
which is also highlighted in yellow in Figure~\ref{fig-xml1}. On sitting
days where the Federation Chamber met as well, a child element exists
for the Federation Chamber too. This allows for the distinction of
proceedings of the Chamber and Federation Chamber. As mentioned in
Section~\ref{sec-description}, each proceeding begins with the business
start. The structure of this business start can be seen between the
nodes highlighted in green in Figure~\ref{fig-xml1}. The content we
parse from the business start are shown highlighted in orange in
Figure~\ref{fig-xml1}. Next comes the debates and nested sub-debates,
which follow the business start. Broadly, the speeches within Hansard
transcripts are structured with a debate node, containing a sub-debate 1
child-node which has a sub-debate 2 child-node nested within it. That
said, sometimes sub-debate 2 is not nested within sub-debate 1. Each of
these three elements as well as their respective sub-elements contain
important information on the title of the debate or sub-debate, who is
speaking, and what is being said. As mentioned, in many cases speeches
are interrupted by interjections, which are embedded in these
sub-elements as well. The final key distinct component of Hansard is
Question Time, in which question and answer elements are classified
differently than general debate text. Sometimes questions and answers
are nested within debate or sub-debate child nodes, and other times they
have their own child node. More detail on the processing of Question
Time will follow in Section~\ref{sec-qa}.

\begin{figure}

{\centering \includegraphics{xml1.png}

}

\caption{\label{fig-xml1}Snapshot of beginning of XML file for Hansard
on 25 February 2020}

\end{figure}

Once code was written to parse all components, it was organized into
individual scripts. The first script includes everything from the
session header element, the second contains everything from the debate
element, and the third and fourth contain everything from the sub-debate
1 and 2 elements, respectively. The fifth script contains debate
interjection information, and the sixth contains content from Question
Time. The final script is a compilation of everything. The next step was
to further develop these scripts to produce tidy data sets from each
parsed element, where each statement is separated onto its own row with
details about the speaker, and rows are placed in chronological order.
This first involved correcting the variable classes and adding a number
of indicator variables to differentiate where statements came from, such
as Chamber versus Federation Chamber or sub-debate 1 versus sub-debate
2. The next key task stemmed from the fact that the raw text data were
not separated by each statement when parsed. In other words, any
interjections, comments made by The Speaker or Deputy Speaker and
continuations within an individual speech were all parsed together as a
single string. As such, the name, name ID, electorate and party details
were only provided for the person who's turn it was to speak. Splitting
up these speeches in a way which would be generalizable across sitting
days required much thought and effort. Section~\ref{sec-interject} will
provide further details on the intricacies of this task.

When beginning to run these scripts on transcripts from other sitting
days, it became clear that not every sitting day contains every possible
XML element. For example, some days did not have sub-debate 2 content,
and some days did not have a Federation Chamber proceeding. To improve
the generalizability of these scripts, if-else statements were embedded
within the code wherever an error might arise due to a missing element.
For example, the entire Federation Chamber block of code is wrapped in
an if-else statement for each script, so that it only executes if what
the code attempts to parse actually exists in the file. Once the script
ran without error for a few recent years of Hansard, we continued to
work backwards until extensive changes in tree structure made our script
incompatible with parsing earlier XML files. Specifically, the earliest
date for which it works is 10 May 2011. Before writing a new script for
parsing earlier Hansard, we decided to prioritize cleaning and
finalizing what we have been able to parse. As such we continued
building our scripts, fixing any issues we noticed in the resulting
datasets, and splitting up any additional sections of the parsed text
onto separate rows where necessary. Specifically, we added a section of
our script to separate out general stage directions. More information on
this separation will be provided in Section~\ref{sec-stage}.

After completing the scripts, they were formatted as functions which
take a file name argument. The final script which parses everything from
Hansard produces one CSV file with the data on all proceedings for each
sitting day (example shown in Figure~\ref{fig-main15rows}).

\hypertarget{sec-qa}{%
\subsection{Question Time}\label{sec-qa}}

A key characteristic of the Australian parliament system is the ability
for the executive government to be held accountable for their decisions.
One core mechanism by which this is achieved is called Question Time.
This is a period on each sitting day in the Chamber where members of the
House can ask ministers two types of questions: questions in writing
which are written in advance, or questions without notice which are
asked verbally in the Chamber and are responded to in real time (House
of Representatives 2021). Parsing the components of Question Time
required a slightly different approach than that of the debate speech,
because of its unique structure in the XML document. Sometimes,
questions in writing are included directly in the
\texttt{"chamber.xscript"} child node, with sub-child nodes called
\texttt{"question"} and \texttt{"answer"} to differentiate the two.
However, in other times, questions in writing are embedded in their own
child node called \texttt{"answers.to.questions"} outside of
\texttt{"chamber.xscript"}. Difficulty arose when trying to order these
questions and answers correctly within the rest of the debate text,
because unlike all other statements made which have both an associated
page number and time stamp, these generally only have a page number.
Time stamps have been a key variable for which we have arranged parsed
text in chronological order, especially since many statements tend to be
on the same page, meaning we cannot rely solely on page number to
capture correct ordering. To approach this issue, we merged all parsed
questions to be in a single dataframe, and all parsed answers in another
dataframe. We then arranged each dataframe by page number, and merged
the two dataframes in such a way that each question would be followed by
its corresponding answer.

\hypertarget{sec-interject}{%
\subsection{Interjections}\label{sec-interject}}

As mentioned, the text was structured and parsed in such a way that
various interjections and comments which happened during a speech were
not separated onto individual rows. Below is a simple example extracted
from Hansard on 30 November 2021, where Mr.~van Manen is interrupted by
The Speaker who states that the time for members' statements has
concluded.

\begin{quote}
``Mr VAN MANEN (Forde---Chief Government Whip) (13:59): It's a great
pleasure to share with the House that Windaroo Valley State High School
has qualified for the finals of the Australian Space Design Competition,
to begin in January next year. The competition is regarded as the
premier STEM competition for high school students and is recognised by
universities around the country. The students are required to respond to
industry-level engineering and requests for tender for design and---The
SPEAKER: Order! In accordance with standing order 43, the time for
members' statements has concluded.''
\end{quote}

As previously discussed, we wanted to have each statement on its own
row, with the correct name, name ID, electorate and party information on
the individual speaking. We approached this task in a number of steps.
Once all parsed text from the XML was merged into one dataframe called
\texttt{main}, we added a \texttt{"speech\_no"} variable. As noted in
Section~\ref{sec-description}, this was done to keep track of which
speech each interjection, comment, or continuation belonged to as we
separated these components onto their own rows.

The next step was to extract all the names and titles preceding these
interjections, comments and continuations. This would enable us to then
separate the speeches in the correct places using these names and titles
in combination with regular expressions, which are patterns of
characters that can be used to search bodies of text. We completed this
extraction process with a few intermediate steps, due to the large
number of name styles and interjection types that had to be accounted
for, each requiring their own unique regular expression format.

In many of the XML files, below the full speech including all
interjections and continuations, there is a series of child-nodes which
exist to capture the structure of interruptions in that speech.
Figure~\ref{fig-xml3} provides an example of this, where the speech was
interrupted by a comment from The Deputy Speaker, and then the member
continued their speech. Looking at the elements highlighted in blue, it
is clear that these child-nodes do not contain the actual text for the
interjection or continuation. As mentioned, this text is embedded within
the speech above it. However, as shown by the content highlighted in
pink in Figure~\ref{fig-xml3}, we were able to extract useful details on
the individual interjecting which we could use later on. Making use of
this structure, we extracted names and information of all individuals
that were categorized within the XML as interjections. We stored this as
a dataframe called \texttt{"interject"}.\footnote{We decided not to
  include this data in our final database, as it is all embedded in our
  resulting datasets which have a flag for interjections.}

\begin{figure}

{\centering \includegraphics[width=3.59375in,height=\textheight]{xml3.png}

}

\caption{\label{fig-xml3}Snapshot of XML structure with interjection and
continuation from 03 February 2021 Hansard}

\end{figure}

We then created lists using both the \texttt{interject} and
\texttt{main} dataframes to capture all the names of individuals who
spoke that day. We added the names of all Members in a number of unique
formats, due to the frequent variation in how names are transcribed in
Hansard. When a Member interjects or continues a speech, the usual form
of their name is a prefix followed by their first name or first initial
and/or last name. There is also variation in the capitalization of these
names.\footnote{Sometimes when someones first name is included, only
  their last name is capitalized, while sometimes their full name is
  capitalized, or other times neither are capitalized.} Another source
of variation is in individuals with more than one first name, as
sometimes only their first first name is written, while other times
their entire first name is written. Additionally, some surnames have
punctuation, and some surnames have specific capitalization such as
``McCormack'', where even in full capitalization, the first ``c''
remains lower case. This variation demands careful consideration when
writing regular expression patterns to search for in the text. In these
lists we also accounted for any general interjection statements
transcribed that were not attributed to an individual, such as ``An
opposition member interjecting-''.

Having these lists enabled us to extract the names of Members and their
associated prefixes as they exist in the text, by searching for exact
matches with regular expression patterns. We then used these extracted
names to split all the speeches up, using a number of regular
expressions with lookarounds. A lookaround can be added to a regular
expression pattern to enhance the specificity of matches. These were
used to ensure that the text was not being split in the wrong places,
such as places where Members were simply being named in the statement of
another Member.

Once all interjections, comments and continuations were successfully
split onto their own rows, we added a general order variable to the
dataset based on row number, to keep track of the order in which
everything was said. The next step was to fill the name, name ID,
electorate and party variables with the correct data for each row. To do
so, we created a lookup table, which contained the unique incomplete
form in which the name was transcribed, and the corresponding full name,
name ID, electorate, and party for that individual
Figure~\ref{fig-lookup}. We used the main dataset from the
\texttt{AustralianPoliticians} package in the creation of this lookup
table (Alexander and Hodgetts 2021).

\begin{figure}

{\centering \includegraphics[width=4.66667in,height=\textheight]{lookup_ex.png}

}

\caption{\label{fig-lookup}First 10 rows of lookup table from 10 May
2018 Hansard processing}

\end{figure}

Next, we merged our main dataframe with the lookup table to replace any
incomplete names with their full names, and to fill in any gaps in the
table with available name ID, electorate and party information. Finally,
we were able to add a flag for interjections. Grouping our data by the
speech number, we defined an interjection as a statement made by anyone
who is not The Speaker, The Deputy Speaker, or the Member whose turn it
was to speak. Figure~\ref{fig-interject} provides an example of a
Federation Chamber proceeding with interjections. Evidently, statements
made by the original speaker Stuart Robert or by The Deputy Speaker
Maria Vamvakinou are not flagged as interjections.

\begin{figure}

{\centering \includegraphics{interject_ex.png}

}

\caption{\label{fig-interject}Example of speech with interjections from
21 November 2016 Hansard}

\end{figure}

\hypertarget{sec-stage}{%
\subsection{Stage Directions}\label{sec-stage}}

As mentioned, one of the final components added to our script was to
separate general stage directions out from statements made by members.
Stage directions are general statements included in the transcript to
document happenings in parliament. Examples of stage directions are
``Bill read a second time'', ``Question agreed to'' or ``Debate
adjourned''. It was unclear to us from the XML and PDF who exactly these
statements were attributed to. For further clarification, we watched
portions of the video recording for some sitting days, and noticed that
where these statements are documented in Hansard, they are not
explicitly stated in parliament. For example, when The Deputy Speaker
says ``The question is that the bill be now read a second time'',
members of the House take a vote, and if the majority is in favour, they
proceed reading the bill the second time. This vote and second reading
is not explicitly transcribed, rather what is written is: ``Question
agreed to. Bill read a second time''. For this reason, we filled the
name variable for these statements with ``stage direction''. Note that
stage directions were not flagged as interjections. These stage
directions are not defined differently from the regular debate speech in
the XML, meaning we had to manually create a list of stage directions to
separate out of the speeches. We have been building upon this list of
stage directions as we work backwards in parsing Hansard.

\hypertarget{sec-conclusion}{%
\section{Conclusion}\label{sec-conclusion}}

Using a number of computational tools, we have created a novel database
for Australian Hansard entirely based on XML data. Despite the public
availability of Hansard, it is not functional in its current form for
analysis at scale. Our work will enable individuals from various
disciplines to perform broad spectrum research on Australian
parliamentary proceedings without requiring a strong computational
background.

Our database provides a CSV file for each sitting day of Hansard from 10
May 2011 to 08 September 2022. In our work, we have embedded a number of
additional variables to the resulting datasets and accounted for
variation in the way that proceedings were initially transcribed. This
ensured that users of our database can easily filter through transcripts
as needed to achieve their distinct research goal.

The vastness of our database and the significance of the content within
it enables the exploration of an endless variety of research questions.
For example, one might be interested in investigating the frequency of
certain words or phrases, such as those relating to Covid-19. Another
area of interest would be frequency of interjections, and the gender of
individuals both being interrupted and doing the interrupting. This
database would also act as a valuable resource for research using
Natural Language Processing tools.

\hypertarget{references}{%
\section*{References}\label{references}}
\addcontentsline{toc}{section}{References}

\hypertarget{refs}{}
\begin{CSLReferences}{1}{0}
\leavevmode\vadjust pre{\hypertarget{ref-alexander2021}{}}%
Alexander, Rohan, and Monica Alexander. 2021. {``The Increased Effect of
Elections and Changing Prime Ministers on Topics Discussed in the
Australian Federal Parliament Between 1901 and 2018.''} \emph{arXiv
Preprint arXiv:2111.09299}.

\leavevmode\vadjust pre{\hypertarget{ref-auspol}{}}%
Alexander, Rohan, and Paul A. Hodgetts. 2021.
\emph{AustralianPoliticians: Provides Datasets about Australian
Politicians}.
\url{https://CRAN.R-project.org/package=AustralianPoliticians}.

\leavevmode\vadjust pre{\hypertarget{ref-lipad}{}}%
Beelen, Kaspar, Timothy Alberdingk Thijm, Christopher Cochrane, Kees
Halvemaan, Graeme Hirst, Michael Kimmins, Sander Lijbrink, et al. 2017.
{``Digitization of the Canadian Parliamentary Debates.''} \emph{Canadian
Journal of Political Science/Revue Canadienne de Science Politique} 50
(3): 849--64.

\leavevmode\vadjust pre{\hypertarget{ref-standingorders}{}}%
Department of the House of Representatives. 2019. \emph{Standing Orders
as at 19 Septmeber 2019}. Commonwealth of Australia.

\leavevmode\vadjust pre{\hypertarget{ref-fraussen2018}{}}%
Fraussen, Bert, Timothy Graham, and Darren R Halpin. 2018. {``Assessing
the Prominence of Interest Groups in Parliament: A Supervised Machine
Learning Approach.''} \emph{The Journal of Legislative Studies} 24 (4):
450--74.

\leavevmode\vadjust pre{\hypertarget{ref-australia2018house}{}}%
House of Representatives. 2018. Edited by D. R. Elder and P. E. Fowler.
\emph{{House of Representatives Practice}}. 7th ed. {Australian
Government - Department of the House of Representatives}.

\leavevmode\vadjust pre{\hypertarget{ref-house2021}{}}%
---------. 2021. \emph{{A window on the house: Practices and procedures
relating to question time}}. {Parliament of Australia}.

\leavevmode\vadjust pre{\hypertarget{ref-Leslie_2021_HoR_Aus}{}}%
Leslie, Patrick. 2021. \emph{Members of the Australian House of
Representatives (1945-2019)}.
\url{https://github.com/palesl/AustralianHouseOfRepresentatives}.

\leavevmode\vadjust pre{\hypertarget{ref-rasiah2010framework}{}}%
Rasiah, Parameswary. 2010. {``A Framework for the Systematic Analysis of
Evasion in Parliamentary Discourse.''} \emph{Journal of Pragmatics} 42
(3): 664--80.

\leavevmode\vadjust pre{\hypertarget{ref-salisbury2011}{}}%
Salisbury, Christopher. 2011. {``'Mr Speaker, i Withdraw...': Standards
of (Mis) Behaviour in the Queensland, Western Australian and
Commonwealth Parliaments Compared via Online Hansard.''}
\emph{Australasian Parliamentary Review} 26 (1): 166--77.

\leavevmode\vadjust pre{\hypertarget{ref-sherratt}{}}%
Sherratt, Tim. 2016. {``Documentation: Historic Hansard.''}
\emph{Historic Hansard}.
\url{http://timsherratt.org/digital-heritage-handbook/docs/historic-hansard/\#}.

\leavevmode\vadjust pre{\hypertarget{ref-vice2017history}{}}%
Vice, John, and Stephen Farrell. 2017. \emph{The History of Hansard}.
House of Lords Library; House of Lords Hansard.

\end{CSLReferences}



\end{document}
